\section{Taxonomy}
\label{sec:taxonomy}

%% ~0.6 page. Table 1 = compact category table. Full statute quotes → appendix / repo.
Our taxonomy has two tiers and two kinds. \textbf{Tier L (Legal)} contains only categories \emph{enumerated} by Korean statute: the Personal Information Protection Act (PIPA) \S23--24 and Enforcement Decree \S18--19 (sensitive and unique identifying information); the Credit Information Act (CIA) \S2 and its Enforcement Decree \S2, which lists 성명·주소·전화번호, e-mail and SNS addresses, the four national identification numbers, \emph{CI/DI and any identifier a credit-information company assigns to a customer} (\S2\,③), and business/corporate registration numbers (\S2\,④); the Real Name Financial Transactions Act \S4 (transaction information); and the Electronic Financial Transactions Act \S2(10) (access media) and Decree \S7\,④ (account and policy numbers, transaction records). \textbf{Tier I (Identifiability)} contains categories no statute enumerates but which act as quasi-identifiers~\cite{pilan2022tab} or, when aggregated, enable profiling and personalised social engineering; membership in Tier I is an empirical judgment, not a legal claim.

Each category is an \textbf{identifier} (a span that alone points at a person or account) or an \textbf{attribute}. CIA \S2(1)(a) defines identifying information as credit information \emph{only when combined with} transaction, creditworthiness or capacity information; we encode this as a span policy: L-identifiers and sensitive information are \nolinkurl{must_mask}, L-attributes are \nolinkurl{mask_if_linked}, and Tier I is \nolinkurl{detect_only}. Table~\ref{tab:taxonomy} lists the 36 categories (17 L-identifier, 7 L-attribute, 12 I-attribute). Staff names are annotated as PII with a \nolinkurl{subject_role} of \emph{staff}; institution and product names, biometric templates and unattributed amounts are excluded.

\begin{table}[t]
\caption{Taxonomy overview: all 36 categories, with abbreviated labels. L denotes Legal and I denotes Identifiability.}
\label{tab:taxonomy}
\small
\begin{tabularx}{\columnwidth}{@{}>{\RaggedRight\arraybackslash}p{0.23\columnwidth}>{\RaggedRight\arraybackslash}X@{}}
\toprule
\textbf{Tier / kind} & \textbf{Categories} \\
\midrule
L / identifier\newline (17) & name; RRN; foreigner reg.\ no.; passport; driver's licence; phone; address; e-mail; online handle; customer ID; business reg.\ no.; corporate reg.\ no.; bank account; securities account; card; contract no.; access credential. \\
\addlinespace
L / attribute\newline (7) & credit transaction; transaction record; delinquency; financial capacity; credit score; public record; sensitive information. \\
\addlinespace
I / attribute\newline (12) & DOB/age; gender; nationality; family relation; employer/position; education; life event; investment profile; consultation content; service usage; device/IP; location mention. \\
\bottomrule
\end{tabularx}
\par\smallskip
\begin{minipage}{\columnwidth}
\RaggedRight\small
\textit{Policies.} L identifiers and sensitive information: must mask.
Other L attributes: mask if linked. I attributes: detect only.
Full subtypes, statutory citations and format rules are in the taxonomy specification.
\end{minipage}
\end{table}
